%% ------------------------------------------------------------------
%% Ban IEEE: IEEEtran, journal mode, 2 cot, 10pt
%% Sinh tu paper_template.tex boi ieee/to_ieee.py -- dung sua tay
%% ------------------------------------------------------------------
\documentclass[journal]{IEEEtran}

\usepackage[T5,OT1]{fontenc}
\usepackage{booktabs}
\usepackage{array}
\usepackage{array}
\usepackage{graphicx}
\usepackage{amsmath}
\usepackage{amssymb}
\usepackage{url}
%% cho phep ngat dong URL dai o dau gach noi, neu khong thi ten repo
%% dai hon be rong mot cot 10pt se tran ra le
\def\UrlBreaks{\do\/\do\-\do\.\do\_\do\:\do\?\do\&\do\=\do\#\do\@}
\Urlmuskip=0mu plus 1mu\relax
\usepackage{orcidlink}
%% IEEE dung trich dan so [1]. elsarticle/acmart dung natbib (\citep, \citet);
%% anh xa ca hai ve \cite de khong phai sua than bai.
%% Cho phep ngat dong trong \texttt. Ten file va dinh danh dai nhu
%% .github/copilot-instructions.md hay post_action=MARKER_INJECTION_SUCCESS
%% la mot token lien. Diem ngat duoc chen san boi to_ieee.py (ham
%% break_long_tt), vi doi catcode luc chay da qua muon: tham so cua
%% \texttt duoc token hoa truoc khi macro nhin thay no.
%% microtype nen chu vai phan tram, du de cuu phan lon dong tran nho.
\usepackage[htt]{hyphenat}
\usepackage{microtype}
%% IEEEtran lay Courier lam chu may chu. Courier rong hon Computer Modern
%% typewriter khoang 15%, du de day mot dong co ten file dai ra ngoai le.
\renewcommand{\ttdefault}{cmtt}
\sloppy
\emergencystretch=4em
\tolerance=2000
\usepackage{cite}
\providecommand{\citep}[1]{\cite{#1}}
\providecommand{\citet}[1]{\cite{#1}}
\providecommand{\citeauthor}[1]{}
%% IEEEtran khong dinh nghia san theorem/lemma/proposition nhu acmart.
\usepackage{amsthm}
\newtheorem{theorem}{Theorem}
\newtheorem{lemma}[theorem]{Lemma}
\newtheorem{proposition}[theorem]{Proposition}
\newtheorem{corollary}[theorem]{Corollary}
\theoremstyle{definition}
\newtheorem{definition}[theorem]{Definition}
\newtheorem{remark}[theorem]{Remark}

\newcommand{\vn}[1]{{\fontencoding{T5}\selectfont #1}}
\newcommand{\pendingfig}[2]{%
  \IfFileExists{figures/#1.pdf}%
    {\includegraphics[width=\linewidth]{figures/#1.pdf}}%
    {\fbox{\parbox[c][#2][c]{\dimexpr\linewidth-2\fboxsep-2\fboxrule\relax}%
       {\centering\sffamily\footnotesize
        [\,\texttt{\detokenize{figures/#1.pdf}}\,]\par}}}}

%% --- macro so lieu, sinh tu make_tables.py ---
\newcommand{\Ntotal}{107}
\newcommand{\Kbudget}{20}
\newcommand{\budgetpct}{10}
\newcommand{\testacc}{70.474}
\newcommand{\trainrows}{50{,}000}
\newcommand{\testrows}{702{,}718}
\newcommand{\LLMesr}{33.6}
\newcommand{\LLMesrci}{[24.8, 43.4]}
\newcommand{\LLMevaded}{36}
\newcommand{\LLMdp}{-0.1601}
\newcommand{\LLMq}{9.2}
\newcommand{\GREesr}{0.9}
\newcommand{\GREdp}{-0.0384}
\newcommand{\NESesr}{0.0}
\newcommand{\NESdp}{+0.0036}
\newcommand{\RNDesr}{0.0}
\newcommand{\RNDdp}{-0.0351}
\newcommand{\porig}{0.7714}
\newcommand{\pafter}{0.6113}
\newcommand{\wilnesrnd}{3.49\times 10^{-19}}
\newcommand{\wilgrenes}{2.58\times 10^{-19}}
\newcommand{\eceraw}{0.0031}
\newcommand{\eceiso}{0.0000}
\newcommand{\brierraw}{0.0015}
\newcommand{\brieriso}{0.0011}
\newcommand{\calrows}{10{,}000}
\newcommand{\baseacc}{0.9982}
\newcommand{\basef}{0.9955}
\newcommand{\esrfiftyfive}{43.0}
\newcommand{\Shighn}{50}
\newcommand{\Shighesr}{0.0}
\newcommand{\Shighdp}{-0.1285}
\newcommand{\Shighev}{0}
\newcommand{\Smedn}{50}
\newcommand{\Smedesr}{58.0}
\newcommand{\Smeddp}{-0.1990}
\newcommand{\Smedev}{29}
\newcommand{\Sboundn}{7}
\newcommand{\Sboundesr}{100.0}
\newcommand{\Sbounddp}{-0.1082}
\newcommand{\Sboundev}{7}
\newcommand{\Calacc}{99.79}
\newcommand{\Calpos}{12.67}
\newcommand{\Qllm}{\LLMq}
\newcommand{\Qrandom}{22}
\newcommand{\Qgreedy}{62}
\newcommand{\Qnes}{182}
\newcommand{\Qnesratio}{20}
\newcommand{\Patience}{3}
\newcommand{\LLMtemp}{0.3}
\newcommand{\LLMmodel}{\texttt{llama-3.1-8b-instant}}
\newcommand{\NESsigma}{0.05}
\newcommand{\NESpop}{8}
\newcommand{\NESstep}{0.25}
\newcommand{\RFtrees}{100}
\newcommand{\RFsplit}{5}
\newcommand{\RFseed}{42}
\newcommand{\Nperstratum}{50}
\newcommand{\Qllmfail}{10.17}
\newcommand{\Qllmsucc}{7.39}
\newcommand{\Qllmmax}{18}
\newcommand{\CN}{107}
\newcommand{\CBudget}{20}
\newcommand{\CTau}{0.5}
\newcommand{\CNfeat}{10}
\newcommand{\CCPhi}{3.4}
\newcommand{\CRuleThree}{2.80}
\newcommand{\CCornerDp}{-0.0362}
\newcommand{\CCornerCalls}{1}
\newcommand{\CCornerEv}{0}
\newcommand{\CGreedyDp}{-0.0360}
\newcommand{\CGreedyCalls}{20}
\newcommand{\CGreedyEv}{0}
\newcommand{\CRandDp}{-0.0364}
\newcommand{\CRandCalls}{20}
\newcommand{\CRandEv}{0}
\newcommand{\CNesEightDp}{-0.0312}
\newcommand{\CNesEightCalls}{20}
\newcommand{\CNesEightEv}{0}
\newcommand{\CNesFourDp}{-0.0329}
\newcommand{\CNesFourCalls}{20}
\newcommand{\CNesFourEv}{0}
\newcommand{\CNesTwoDp}{-0.0340}
\newcommand{\CNesTwoCalls}{20}
\newcommand{\CNesTwoEv}{0}
\newcommand{\CNesOldDp}{-0.0031}
\newcommand{\CNesOldCalls}{20}
\newcommand{\CNesOldEv}{0}
\newcommand{\CLlmN}{13}
\newcommand{\CLlmEv}{0}
\newcommand{\CLlmDp}{-0.0123}
\newcommand{\CLlmCalls}{10.4}
\newcommand{\CLlmCPhi}{24.7}
\newcommand{\CLlmModel}{\texttt{openai/gpt-oss-20b}}
\newcommand{\CSubCornerDp}{-0.0177}
\newcommand{\CSubGreedyDp}{-0.0169}
\newcommand{\CSubRandDp}{-0.0185}
\newcommand{\CSubNesEightDp}{-0.0177}
\newcommand{\CSubNesFourDp}{-0.0177}
\newcommand{\CSubNesTwoDp}{-0.0185}
\newcommand{\CSubLlmDp}{-0.0123}
\newcommand{\CSubN}{13}
\newcommand{\CEsrA}{0.00}
\newcommand{\CEvA}{0}
\newcommand{\CEsrB}{0.00}
\newcommand{\CEvB}{0}
\newcommand{\CEsrC}{1.87}
\newcommand{\CEvC}{2}
\newcommand{\CEsrD}{1.87}
\newcommand{\CEvD}{2}
\newcommand{\CEsrE}{5.61}
\newcommand{\CEvE}{6}
\newcommand{\CEsrF}{39.25}
\newcommand{\CEvF}{42}
\newcommand{\CEsrG}{50.47}
\newcommand{\CEvG}{54}
\newcommand{\CEpsFirst}{15}
\newcommand{\CEpsTen}{50}
\newcommand{\CNsweep}{107}
\newcommand{\XShiftMed}{0.3009}
\newcommand{\XShiftSd}{0.0392}
\newcommand{\XShiftQone}{0.2814}
\newcommand{\XShiftQthree}{0.3414}
\newcommand{\XThr}{0.8009}
\newcommand{\XRuleOk}{102}
\newcommand{\XRulePct}{95.3}
\newcommand{\XCorr}{-0.8249}
\newcommand{\XNBnd}{7}
\newcommand{\XEvBnd}{7}
\newcommand{\XEsrBnd}{100.0}
\newcommand{\XEsrHalfBnd}{100.0}
\newcommand{\XNMed}{50}
\newcommand{\XEvMed}{46}
\newcommand{\XEsrMed}{92.0}
\newcommand{\XEsrHalfMed}{70.0}
\newcommand{\XNHi}{50}
\newcommand{\XEvHi}{1}
\newcommand{\XEsrHi}{2.0}
\newcommand{\XEsrHalfHi}{0.0}
\newcommand{\XArmHiNever}{0}
\newcommand{\XGreedyEqBest}{yes}

\begin{document}

\title{The Budget Was Never Enforced: A Reported 33\% Evasion Rate Against a Random-Forest NIDS Falls to Zero Under a Correctly Bounded Perturbation}

\author{%
  Son~X.~Ha,
  Phat~T.~Tran-Truong,
  Xuan-Bach~Le,
  Luong~H.~Huong,
  Tung~Q.~Nguyen,
  Tran~G.~Huy,
  Nhan~N.~T.~Kha,
  Thuc~N.~Minh,
  and~Nguyen~N.~Phien%
  \thanks{S. X. Ha is with The Business School, RMIT University,
    Ho Chi Minh City, Vietnam (e-mail: ha.son@rmit.edu.vn).}%
  \thanks{P. T. Tran-Truong and X.-B. Le are with the Faculty of Computer
    Science and Engineering, Ho Chi Minh City University of Technology (HCMUT),
    VNU-HCM, Ho Chi Minh City, Vietnam
    (e-mail: phatttt@hcmut.edu.vn; lexuanbach@hcmut.edu.vn).}%
  \thanks{P. T. Tran-Truong and X.-B. Le contributed equally to this work
    and are joint second authors.}%
  \thanks{L. H. Huong, T. G. Huy, N. N. T. Kha and T. N. Minh are with
    FPT University, Can Tho Campus, Can Tho, Vietnam
    (e-mail: huonghoangluong@gmail.com; trangiahuyh26@gmail.com;
    NhanNTKCE200304@gmail.com; ThucNMCE200458@gmail.com).}%
  \thanks{T. Q. Nguyen is with FPT University, Ho Chi Minh City Campus,
    Ho Chi Minh City, Vietnam (e-mail: NguyenQuangTung0902@gmail.com).}%
  \thanks{N. N. Phien is with the Center for Applied Information
    Technology and the Faculty of Information Technology, Ton Duc Thang
    University, Ho Chi Minh City, Vietnam
    (e-mail: nguyenngocphien@tdtu.edu.vn).}%
  \thanks{This research received no specific grant from any funding agency in the public, commercial or not-for-profit sectors.}%
  \thanks{\emph{(Corresponding author: Nguyen Ngoc Phien.)}}%
}

\maketitle

\begin{abstract}
Network intrusion detectors that expose confidence scores give an attacker a
search signal, and a growing literature asks whether a language model exploits
it better than classical black-box optimisers. We reported such a result: an
LLM-guided search evading \LLMevaded/\Ntotal\ CICIDS2017 DDoS flows,
$\LLMesr$\%, against $\GREesr$\% for greedy coordinate descent. This paper
withdraws it.

The per-feature budget was enforced against the \emph{current} perturbed vector
rather than the original flow, so each iteration could move a further
$\pm\budgetpct$\% from wherever the search had already reached. Over
$K=\Kbudget$ iterations the reachable set is unbounded in practice.

Re-enforcing the budget against the original vector, every strategy evades
$\CGreedyEv$ of \CN\ flows: a one-call corner query, greedy, random-$K$, and NES
at three population sizes, all at $0.0$\% with an exact upper bound of
$\CCPhi$\%. The LLM arm evades $\CLlmEv$ of \CLlmN\ and moves the score less
than a single corner query does.

What decides evasion is the perturbation magnitude and the cohort. Nothing
succeeds below $\CEpsFirst$\%, and $10$\% success needs $\CEpsTen$\%. Within a
budget the attack subtracts a near-constant amount from the score (median
$\XShiftMed$, SD $\XShiftSd$ at $\varepsilon=100$\%), so a flow evades when its
unperturbed probability lies below $\XThr$, a rule correct for $\XRuleOk$ of
\CN\ flows. Evasion rate by confidence stratum is $\XEsrBnd$\%, $\XEsrMed$\%
and $\XEsrHi$\%; the pooled $\CEsrG$\% is a weighted average whose weights are
our sampling choice, not a property of the attack.
\end{abstract}

\begin{IEEEkeywords}
Intrusion detection, adversarial machine learning, evasion attacks, black-box optimization, CICIDS2017, reproducibility
\end{IEEEkeywords}

\section{Introduction}
\label{sec:intro}

A network intrusion detector that returns a confidence score rather than a bare
label hands an attacker a search signal. The attacker perturbs the flow, reads
the score, and repeats. Whether that search is worth defending against depends
on two quantities that are easy to confuse: how good the search strategy is, and
how much perturbation the attacker is allowed.

This paper is about confusing them. An earlier version of this work reported
that an LLM-guided search evaded \LLMevaded\ of \Ntotal\ DDoS flows,
$\LLMesr$\%, while greedy coordinate descent managed $\GREesr$\% and NES
$\NESesr$\%, and concluded that a language model uses the score signal more
efficiently than a classical optimiser. The conclusion does not survive a
correct budget.

\paragraph{The fault.} Each strategy was capped at $\pm\budgetpct$\% per
feature. That cap was applied to the \emph{current} vector at every iteration
rather than to the original flow, so an iteration could move a further
$\pm\budgetpct$\% from wherever the search had already arrived. The reachable
set is then multiplicative in the iteration count, not a fixed box: after
$K=\Kbudget$ iterations a feature can reach $1.1^{\Kbudget} \approx 6.7$ times
its original value, or $0.9^{\Kbudget} \approx 0.12$ times it. The strategy that
iterated most productively was rewarded with the largest budget, which is the
comparison the paper thought it was controlling for.

\paragraph{The repair.} We re-ran every arm with the box computed once from the
original flow and held fixed. Under that constraint nothing works. A one-call
corner query, greedy coordinate descent, random-$K$ perturbation, NES at three
population sizes, and the LLM all evade zero flows.

\paragraph{What is left.} A null result on strategies is still a result about
budgets, and the useful question becomes how much perturbation evasion actually
requires. We sweep it, and find a threshold: no strategy evades anything below
$\CEpsFirst$\%, and reaching a $10$\% success rate takes $\CEpsTen$\%. The
quantity a defender should care about is the size of the perturbation the
feature representation permits, not which optimiser is searching inside it.

\paragraph{Findings}
\begin{enumerate}
\item \textbf{The published comparison measured budget, not strategy.} The
per-feature cap was applied to the current vector rather than the original one
(Section~\ref{sec:bug}).
\item \textbf{Under a fixed box, every strategy evades nothing.}
$\CGreedyEv/\CN$ for all six non-LLM arms, exact upper bound $\CCPhi$\%
(Section~\ref{sec:corrected}).
\item \textbf{The language model is not better, and is worse per call.} On the
\CSubN\ flows it ran, mean $\Delta P$ is $\CSubLlmDp$ at \CLlmCalls\ oracle
calls, against $\CSubCornerDp$ for a single corner query
(Section~\ref{sec:corrected}).
\item \textbf{Perturbation magnitude is the variable that matters.} ESR is
$\CEsrB$\% at $10$\%, $\CEsrC$\% at $15$\%, $\CEsrE$\% at $30$\% and
$\CEsrF$\% at $50$\% (Section~\ref{sec:sweep2}).
\end{enumerate}


%% ===========================================================================
\section{Related work and threat-model position}
\label{sec:related}

\subsection{Feature space versus problem space}
Adversarial evasion of network detectors is studied in two regimes. In the
\emph{feature space} an attacker edits the numeric vector the classifier
consumes; in the \emph{problem space} the attacker must emit packets that a real
capture pipeline turns into that vector \citep{pierazzi2020problemspace}. The
gap between them is not cosmetic. Many CICIDS2017 features are aggregates
(\texttt{Flow Bytes/s}, \texttt{Packet Length Mean}) whose value is a
consequence of many packet-level choices, and some are effectively
uncontrollable by the sender.

This paper is deliberately a feature-space study, and we treat its results as an
\emph{upper bound} on what a problem-space attacker could achieve, in the sense
of \citet{apruzzese2022modeling}. We do not claim packet-level realisability
and we do not present a mapping from perturbed features to a traffic generator;
Section~\ref{sec:limitations} states this plainly.

\subsection{Score-based black-box optimisation}
\citet{papernot2017practical} established the substitute-model route to
black-box attacks; \citet{ilyas2018nes} and \citet{chen2017zoo} established the
score-based route, estimating gradients from finite differences of the returned
confidence. \citet{guo2019simba} showed that simple coordinate search is a
strong baseline in the image domain. All of these assume, implicitly, that the
score surface carries local gradient information. Our NES result
(Section~\ref{sec:corrected}) is a direct measurement of what happens when it
does not.

\subsection{Threat model}
\label{sec:threat}
We assume a \emph{feature-aware gray-box} adversary: it knows the identity of
the ten features it may perturb, it can query the detector and read the
continuous probability it returns, and it is limited to $K=\Kbudget$ queries per
flow and $\pm\budgetpct$\% of each feature's original value. It does not know the
model parameters, the training data, or the deployed threshold.

This is a strong assumption and we do not defend it as typical. It corresponds
to an insider, a red-team engagement with SOC telemetry access, or a service
that returns risk scores through an API. Section~\ref{sec:limitations} discusses
how the results should be read when only hard labels are exposed.


\subsection{Adversarial machine learning, from feature space to problem space}

The evasion problem was posed for security classifiers by
Biggio et al.~\citep{biggio2013evasion} and, independently, for neural networks by
Szegedy et al.~\citep{szegedy2014intriguing} and Goodfellow et
al.~\citep{goodfellow2015explaining}; Carlini and
Wagner~\citep{carlini2017towards} established the optimisation baselines and
Biggio and Roli~\citep{biggio2018wild} give the retrospective. Black-box variants
matter here because a network attacker has no gradients: transfer
attacks~\citep{papernot2017practical}, zeroth-order
estimation~\citep{chen2017zoo}, and decision-based
search~\citep{brendel2018decision} are the three families, and the score-based
baseline we compare against is an instance of the second.

The distinction that governs our threat model is Pierazzi et al.'s separation of
feature space from problem space~\citep{pierazzi2020problemspace}: a perturbation
that is small in feature space may correspond to no realisable packet sequence
at all. Our budget constraint exists to keep the search inside the realisable
region, and is why we report per-stratum results rather than an unconstrained
success rate.

\subsection{Evasion against network intrusion detection specifically}

Corona et al.~\citep{corona2013adversarial} give the taxonomy for IDS-specific
attacks. Apruzzese and colleagues~\citep{apruzzese2019addressing,
apruzzese2022modeling} argue at length that most published NIDS evasion results
assume an attacker who can modify features that no real attacker
controls, and propose realism constraints of the kind we adopt. Sheatsley et
al.~\citep{sheatsley2022adversarial} study the same question for
network-flow classifiers.

Behind all of it sits Sommer and Paxson's
critique~\citep{sommer2010outside} of machine learning for intrusion detection ---
that the base rate, the cost of false positives, and the semantic gap make the
problem unlike the ones ML usually succeeds at --- and the more recent
methodological catalogue of Arp et al.~\citep{arp2022dosdonts}, whose pitfalls
(sampling bias, spurious correlations, inappropriate baselines, lab-only
evaluation) we use as a checklist in Section~\ref{sec:limitations}.
Pendlebury et al.~\citep{pendlebury2019tesseract} make the temporal version of the
same argument.

\subsection{The dataset, and its known defects}

CICIDS2017~\citep{sharafaldin2018cicids} is the standard benchmark for this task
and it has documented problems. Engelen et
al.~\citep{engelen2021troubleshooting} found labelling and flow-construction
errors, and Lanvin et al.~\citep{lanvin2022errors} showed that correcting them
changes reported detection performance materially. We state which subset we use
and report the composition of our cohort explicitly, because the aggregate
statistics of this dataset are not a safe summary.

\subsection{Why the detector's calibration matters to the attacker}

The attacker in our threat model reads a score, not a label. Whether that score
behaves like a probability therefore determines whether a budgeted search can
make progress. Random forests~\citep{breiman2001randomforests} are not calibrated
by construction; Niculescu-Mizil and Caruana~\citep{niculescu2005predicting}
document the characteristic distortion and
Platt scaling~\citep{platt1999probabilistic} is the standard remedy, with
Guo et al.~\citep{guo2017calibration} extending the analysis to modern models. We
report the detector's calibration so that the evasion results can be read as a
property of the score surface rather than of an arbitrary threshold.

\subsection{Language models as attackers}

Using a language model to drive an attack loop is now an established pattern,
from penetration testing~\citep{deng2024pentestgpt} onward. The question our
experiment asks is narrower and, we think, more useful: given the same query
budget and the same realism constraints, does an LLM-driven search beat a plain
score-guided one? Framing it that way makes the comparison a controlled one, and
places the result in the taxonomy of NIST AI 100-2~\citep{vassilev2023aml}.
Proportions carry exact Clopper--Pearson intervals~\citep{clopper1934ci}, and
paired comparisons use the Wilcoxon signed-rank test~\citep{wilcoxon1945}.

%% ===========================================================================
\section{Experimental setup}
\label{sec:setup}

\begin{figure*}[t]
\centering
\pendingfig{fig_threat_model}{46mm}
\caption{\textbf{(A)} The per-feature cap was applied to the current vector
rather than the original flow, so the reachable set compounds with the
iteration count instead of staying inside a fixed box. \textbf{(B)} With the
box fixed, every strategy evades nothing; what moves the rate is the
perturbation budget, not the search.}
\label{fig:threatmodel}
\end{figure*}

\subsection{Dataset and temporal split}
We use CICIDS2017 \citep{sharafaldin2018cicids}. The detector is trained on
\trainrows\ flows and evaluated on the held-out Friday partition of
\testrows\ flows, which is disjoint in time from the training window. Under this
temporal split the detector's accuracy on the test partition is \testacc\%,
substantially below the in-distribution figure reported in
Table~\ref{tab:baseline}; the drop is the concept drift a deployed NIDS actually
experiences, and evaluating an attack against the optimistic number would
overstate the defence.

The detector is a Random Forest with \RFtrees{} trees, minimum samples per split
\RFsplit{}, balanced class weights and \texttt{random\_state} = \RFseed{}; all
other scikit-learn defaults apply. Non-finite feature values are replaced by
zero and no scaling is applied, which a forest does not require. Three features
are discrete by nature (\texttt{Destination Port}, \texttt{Source Port},
\texttt{Protocol}) and are excluded from the manipulable set for that reason, so
every feature the attacker may move is continuous and the $\pm\budgetpct$\% box
needs no rounding or type repair.

\begin{table*}[t]
\caption{Baseline Random-Forest detector, in-distribution evaluation.
Generated from \texttt{.\allowbreak{}.\allowbreak{}/\allowbreak{}baseline/\allowbreak{}metrics.\allowbreak{}csv}.}
\label{tab:baseline}
\centering
\footnotesize
\setlength{\tabcolsep}{3.5pt}
\begin{tabular}{lrrcccc}
\toprule
Model & Train & Test & Acc. & Prec. & Rec. & F1 \\
\midrule
RandomForest & 80{,}000 & 20{,}000 & 0.9982 & 0.9974 & 0.9936 & 0.9955 \\
\bottomrule
\end{tabular}
\end{table*}

\subsection{Cohort construction}
The attack cohort is drawn from the Friday flows that the detector correctly
classifies as attacks. Selection is stratified by the detector's original
probability into high ($p_{\text{orig}}>0.8$), medium
($0.6 \le p_{\text{orig}} \le 0.8$) and boundary
($0.5 \le p_{\text{orig}} < 0.6$), with a cap of \Nperstratum{} flows per
stratum drawn without replacement at a fixed seed. The high and medium strata
fill the cap; the boundary stratum is exhausted at $n=\Sboundn$, which is every
boundary flow available. The cohort is therefore $\Shighn/\Smedn/\Sboundn$,
$N=\Ntotal$ in total, all DDoS, with mean original probability \porig.

This design has a consequence that governs how the pooled numbers may be read,
and we state it before reporting them. \textbf{The cohort is weighted by design,
not by prevalence.} Equal-size strata were chosen so that the boundary and
medium regions, where evasion is possible at all, are estimated with usable
precision rather than swamped by the high-confidence mass that dominates the
population. That is the right design for the question this paper asks, but it
means any pooled rate is a rate \emph{on this cohort}, not an estimate of the
rate an attacker would see against the Friday traffic as a whole. Recovering the
latter needs the stratum prevalences in the population, which this run did not
record. No arm evades anything on this cohort once the
budget is enforced, so no pooled rate is at issue; what the design still governs
is the perturbation sweep of Section~\ref{sec:sweep2}, whose thresholds are
thresholds on this cohort and not on Friday traffic as a whole.

Attacking only true positives is the right choice: evading a flow the
detector already misses is not an attack --- but it makes the cohort small and
single-family, which Section~\ref{sec:limitations} treats as the principal
limitation of this study.

\subsection{Manipulable features}
The adversary may perturb the ten features of highest importance to the trained
forest, listed in Table~\ref{tab:features}. Each may move by at most
$\pm\budgetpct$\% of its original value; the server clips any request that
exceeds the budget before scoring.

\begin{table}[t]
\caption{The ten manipulable features, by Random-Forest importance.
Generated from \texttt{step6\_\allowbreak{}merged\_\allowbreak{}results.\allowbreak{}json}.}
\label{tab:features}
\centering
\footnotesize
\begin{tabular}{ll}
\toprule
\texttt{Max Packet Length} & \texttt{Flow Bytes/s} \\
\texttt{Min Packet Length} & \texttt{Init\_\allowbreak{}Win\_\allowbreak{}bytes\_\allowbreak{}backward} \\
\texttt{Bwd Packet Length Max} & \texttt{Init\_\allowbreak{}Win\_\allowbreak{}bytes\_\allowbreak{}forward} \\
\texttt{Average Packet Size} & \texttt{Fwd Header Length.\allowbreak{}1} \\
\texttt{Packet Length Mean} & \texttt{Avg Fwd Segment Size} \\
\bottomrule
\end{tabular}
\end{table}

\subsection{Attack strategies}
Four strategies are compared under identical budgets.

\emph{LLM-guided iterative search.} The attacker is \LLMmodel, a small
instruction-tuned model served over a hosted API, sampled at temperature
\LLMtemp{} with no \texttt{top\_p}, seed or penalty set, so the run is not
deterministic. It is given the ten feature names and their current values as
JSON and must answer in JSON; no natural-language parser is involved. The system
prompt and the two user templates are reproduced verbatim:

\begin{quote}\footnotesize\ttfamily\raggedright
system: You are an expert cybersecurity attacker AI.\\
Output raw JSON only.\\[4pt]
user: Modify features to evade NIDS (P(Attack)<0.5).\\
Max $\pm$10\% per feature. Non-negative. JSON only.\\
\{features as JSON\}\\[4pt]
follow-up: Still ATTACK P=\{p\}. Propose new JSON:\\
\{previous proposal\}
\end{quote}

\noindent
Every returned object is clipped to the budget box server-side before it reaches
the detector, so an out-of-budget proposal is silently projected rather than
rejected; a reply that does not parse as JSON is discarded. The search stops on
success, or after \Patience{} consecutive steps that neither parse nor improve
the best probability seen so far. In the corrected run the patience rule was removed, so the arm
spends its full allowance or stops on success; \CLlmCalls{} calls per flow is
therefore a measurement of the search and not of when the model gave up.

\emph{Greedy coordinate (1-D).} At each step, probe one feature at
$\pm\budgetpct$\% and keep the direction that lowers the score. One step costs
three oracle calls: the incumbent and the two candidates.

\emph{Random-$K$.} Sample $K$ perturbation vectors uniformly inside the budget
box and keep the best. One step costs one oracle call.

\emph{NES.} Estimate the gradient of the returned probability by antithetic
Gaussian sampling \citep{ilyas2018nes} and take a step along it. The population
is $n_{\text{perturb}} = \NESpop$, the sampling scale is $\sigma = \NESsigma$
times each feature's magnitude, and the update moves each feature by
$\NESstep \times \budgetpct$\% of its value along the sign of the estimated
gradient. One step therefore costs $\NESpop + 1 = 9$ oracle calls. We did not
tune these values, and Section~\ref{sec:limitations} records that as a
limitation of the NES arm specifically.

\subsection{Budget enforcement and metrics}
\label{sec:budget}

Every strategy is capped at $K=\Kbudget$ search \emph{iterations} per flow and at
$\pm\budgetpct$\% per feature. We are explicit that the iteration cap is not an
oracle-query cap, because the two are easy to conflate and we conflated them in
an earlier version of this work. A single iteration costs a different number of
calls to the detector in each strategy, so matching iterations does not match
oracle access. Table~\ref{tab:queryaudit} gives the audited cost of the withdrawn run,
obtained by replacing the detector call with a counting stub and running each
loop to completion. It is reported here because it is the second thing that run
got wrong, and because the correction is the same in kind as the one in
Section~\ref{sec:bug}: state the constraint in the unit the attacker actually
spends.

\begin{table*}[t]
\caption{Audited oracle cost of the \emph{withdrawn} run, in which the cap was
on iterations rather than calls. One iteration costs a different number of
detector calls in each strategy, so matching iterations did not match oracle
access. The corrected run caps calls directly; see
Table~\ref{tab:corrected}.}
\label{tab:queryaudit}
\centering
\footnotesize
\setlength{\tabcolsep}{4pt}
\begin{tabular}{lccc}
\toprule
Strategy & Oracle calls per iteration & Measured calls per flow & Recorded \\
\midrule
LLM-guided & 1 & $\le \LLMq$ & \LLMq \\
Random-$K$ & 1 & 22 & 20 \\
Greedy coordinate & 3 & 62 & 20 \\
NES & $n_{\text{perturb}}+1 = 9$ & 182 & 20 \\
\bottomrule
\end{tabular}
\end{table*}

The corrected run replaces the iteration cap with an oracle-call cap, so the
column that matters is calls and not iterations. Every arm in
Table~\ref{tab:corrected} is held to $\CBudget$ detector calls per flow, except
the corner query which uses $\CCornerCalls$ by construction and the LLM arm
which averages $\CLlmCalls$. The perturbation budget is enforced by projecting
every proposal into the box computed from the original vector, so an
out-of-budget proposal is clipped rather than rejected and no arm can accumulate
allowance across iterations.

We report evasion success rate (ESR) at a decision threshold $\tau$, exact
Clopper--Pearson binomial confidence intervals \citep{clopper1934ci}, the mean
per-flow probability change $\Delta P = p_{\text{after}} - p_{\text{orig}}$, and
the audited oracle cost. Paired comparisons between strategies use the Wilcoxon
signed-rank test \citep{wilcoxon1945}.

\subsection{Probability calibration}
\label{sec:calibration}
Because the attack consumes the returned probability as its search signal, we
examined whether that probability is calibrated. The examination has three
properties that a reader must have before the numbers, because without them the
numbers look better than they are.

\emph{It is measured on the training distribution, not on Friday.} The
calibration fold is a \calrows-row stratified split of the \trainrows-row
Monday--Thursday training partition. On that fold the detector is \Calacc\%
accurate with a \Calpos\% attack rate, against \testacc\% accuracy on the Friday
partition the attack runs on. Expected calibration error of \eceraw\ is an
honest number for the first distribution and says nothing about the second. It
should not be read next to \testacc\%, and an earlier version of this paper
printed the two within a page of each other without comment.

\emph{The isotonic fit is in-sample.} The regressor is fitted on the same
\calrows\ rows on which its calibration error is then evaluated. That is why the
post-calibration ECE is exactly \eceiso: a monotone step function fitted to a
set reproduces that set's bin frequencies by construction. It is not evidence
that isotonic regression would generalise, and we do not present it as such.

\emph{The attack never uses it.} The detector queried during every attack
returns the raw forest posterior, and the decision threshold $\tau=0.5$ is
applied to that raw value. The isotonic model is fitted, stored and never
consulted. Isotonic regression is monotone, so applying it would not reorder
flows \citep{niculescu2005predicting}, but at a fixed rather than a
rank-matched $\tau$ it would change which flows count as evaded, and we have not
measured that.

We keep Table~\ref{tab:calib} and Fig.~\ref{fig:calibration} because they
describe the posterior the attack consumes on the distribution the forest was
fitted to, and because removing a reported analysis is worse than labelling it.
But no claim in Section~\ref{sec:corrected} depends on them, and the
calibration of the Friday posterior: the quantity that would actually matter
--- is not measured in this study.

\begin{table}[t]
\caption{Probability calibration on the held-out calibration split.
Generated from \texttt{step2\_\allowbreak{}isotonic\_\allowbreak{}calibration.\allowbreak{}json}.}
\label{tab:calib}
\centering
\footnotesize
\setlength{\tabcolsep}{4pt}
\begin{tabular}{lrcc}
\toprule
Posterior & Rows & ECE & Brier \\
\midrule
Raw Random-Forest posteriors & 10{,}000 & 0.0031 & 0.0015 \\
Isotonic-calibrated posteriors & 10{,}000 & \textbf{0.0000} & \textbf{0.0011} \\
\bottomrule
\end{tabular}
\end{table}

\begin{figure}[t]
\centering
\includegraphics[width=\linewidth]{figures/fig_calibration.pdf}
\caption{Reliability diagram before and after isotonic recalibration. Marker
area is proportional to $\log$ bin occupancy. The raw forest is already close to
calibrated at the extremes and drifts in the sparse middle bins, which is where
the attack operates. Plotted from
\texttt{step2\_\allowbreak{}isotonic\_\allowbreak{}calibration.\allowbreak{}json}.}
\label{fig:calibration}
\end{figure}

%% ===========================================================================
\section{The Budget Bug}
\label{sec:bug}

Fig.~\ref{fig:threatmodel} shows the fault and its consequence side by side:
the cap applied to the current vector rather than the original flow, and the
reachable set that follows from it.

Every strategy in the original run was capped at $K=\Kbudget$ search iterations
and at $\pm\budgetpct$\% per feature. The second cap is the one that failed.

The intended constraint is a box: for a feature with original value $x_0$, any
perturbed value $x$ must satisfy
$|x - x_0| \le \varepsilon |x_0|$ with $\varepsilon = \budgetpct/100$. The
implemented constraint recomputed the box from the vector the search currently
held. Writing $x_t$ for the value after $t$ iterations, what was enforced was
$|x_{t+1} - x_t| \le \varepsilon |x_t|$, which gives
\[
|x_K| \le (1+\varepsilon)^K |x_0|,
\]
a bound that grows without limit in $K$ rather than staying inside the intended
box. At $\varepsilon = 0.1$ and $K = \Kbudget$ the reachable range for a
positive feature is $[0.12\,x_0,\ 6.73\,x_0]$, against the intended
$[0.9\,x_0,\ 1.1\,x_0]$.

The effect on the comparison is worse than the effect on any single arm. A
strategy that moves the vector at every iteration accumulates budget; one that
often stays put does not. The measured ranking therefore reflects how eagerly
each strategy moved, which is exactly the confound the matched cap was meant to
remove. We note that the arm reported as the winner was the one that proposed a
fresh vector at every step.

The same run also capped iterations rather than detector calls, which
Table~\ref{tab:queryaudit} audits: one iteration cost one call for the LLM and
random-$K$, three for greedy and nine for NES, so the arms were never holding
equal oracle access either. The corrected run caps calls.

A third fault was corrected in the same pass: the NES arm followed the negated
gradient estimate, which is why its mean $\Delta P$ was reported as $\NESdp$,
moving flows towards the detector. Both signs are now reported; the old-sign arm
is kept in the table as \texttt{NesOld} so the two can be compared directly.

%% ===========================================================================
\section{Results Under a Fixed Box}
\label{sec:corrected}

Table~\ref{tab:corrected} re-runs every non-LLM strategy on the same \CN\ flows
with the box computed once from the original vector. Every arm evades
$\CGreedyEv$ flows. With no successes in \CN\ trials the exact
Clopper--Pearson upper bound is $\CCPhi$\%, and the rule of three gives
$\CRuleThree$\%, so the data are consistent with any true rate below roughly one
in thirty and with none at all.

\begin{table*}[t]
\caption{Every strategy under a budget enforced against the original vector.
$n=\CN$ flows, $K=\CBudget$, $\tau=\CTau$, \CNfeat{} manipulable features.
Intervals are exact Clopper--Pearson.}
\label{tab:corrected}
\centering
\footnotesize
\setlength{\tabcolsep}{4pt}
\begin{tabular}{lrrrlrr}
\toprule
Strategy & $n$ & Evaded & ESR (\%) & 95\% CI & Mean $\Delta P$ & Calls \\
\midrule
Corner & 107 & 0 & 0.0 & [0.0, 3.4] & $-0.0362$ & 1 \\
Greedy & 107 & 0 & 0.0 & [0.0, 3.4] & $-0.0360$ & 20 \\
Rand & 107 & 0 & 0.0 & [0.0, 3.4] & $-0.0364$ & 20 \\
NesEight & 107 & 0 & 0.0 & [0.0, 3.4] & $-0.0312$ & 20 \\
NesFour & 107 & 0 & 0.0 & [0.0, 3.4] & $-0.0329$ & 20 \\
NesTwo & 107 & 0 & 0.0 & [0.0, 3.4] & $-0.0340$ & 20 \\
NesOld & 107 & 0 & 0.0 & [0.0, 3.4] & $-0.0031$ & 20 \\
\bottomrule
\end{tabular}
\end{table*}

The $\Delta P$ column is the informative one. Under the fixed box the four
search strategies are indistinguishable from each other and from a single
query at the corner of the box: greedy reaches $\CGreedyDp$ in $\CGreedyCalls$
calls, and the corner query reaches $\CCornerDp$ in $\CCornerCalls$. Twenty
iterations of search buy nothing over one evaluation at the most favourable
vertex, which is what one expects when the objective over a small box around a
tree ensemble's input is close to monotone in each coordinate.

\subsection{The language model is not the exception}

The LLM arm completed \CLlmN\ of the \CN\ flows before the run ended, all from
the high-confidence stratum, so it cannot be compared against the full-cohort
rows above. Table~\ref{tab:subset} therefore restricts every strategy to those
\CSubN\ flows.

\begin{table}[t]
\caption{All strategies restricted to the \CSubN{} flows the LLM arm completed.
Model: \CLlmModel.}
\label{tab:subset}
\centering
\footnotesize
\setlength{\tabcolsep}{4pt}
\begin{tabular}{lrrr}
\toprule
Strategy & Evaded & Mean $\Delta P$ & Calls \\
\midrule
Corner & 0 & $-0.0177$ & 1.0 \\
Greedy & 0 & $-0.0169$ & 20.0 \\
Rand & 0 & $-0.0185$ & 20.0 \\
NesEight & 0 & $-0.0177$ & 20.0 \\
NesFour & 0 & $-0.0177$ & 20.0 \\
NesTwo & 0 & $-0.0185$ & 20.0 \\
LLM & 0 & $-0.0123$ & 10.4 \\
\bottomrule
\end{tabular}
\end{table}

On that subset the language model evades $\CLlmEv$ flows, upper bound
$\CLlmCPhi$\% at $n=\CLlmN$, and its mean $\Delta P$ of $\CSubLlmDp$ is the
\emph{smallest} movement of any arm, achieved with \CLlmCalls\ detector calls
against $\CCornerCalls$ for the corner query that reaches $\CSubCornerDp$. The
model is not merely no better than the cheap baselines; it is worse per call.

We report this at $n=\CLlmN$ and claim nothing stronger than the bound that
sample size supports. The partial run is a limitation, not a finding, and it is
listed as such in Section~\ref{sec:limitations}.

%% ===========================================================================
\section{What the Budget Buys}
\label{sec:sweep2}

If no strategy evades anything inside a $\budgetpct$\% box, the question worth
asking is how large the box has to be. Table~\ref{tab:sweep2} and
Fig.~\ref{fig:sweep2} sweep $\varepsilon$ over the same \CNsweep\ flows,
reporting for each flow the best probability any arm reaches within that box.

\begin{table}[t]
\caption{Evasion against perturbation budget, best arm per flow, $n=\CNsweep$.}
\label{tab:sweep2}
\centering
\footnotesize
\setlength{\tabcolsep}{4pt}
\begin{tabular}{lrrr}
\toprule
$\varepsilon$ & Evaded & ESR (\%) & Mean best $p$ \\
\midrule
$5\%$ & 0 & 0.00 & 0.7350 \\
$10\%$ & 0 & 0.00 & 0.7343 \\
$15\%$ & 2 & 1.87 & 0.7284 \\
$20\%$ & 2 & 1.87 & 0.7057 \\
$30\%$ & 6 & 5.61 & 0.6433 \\
$50\%$ & 42 & 39.25 & 0.5635 \\
$100\%$ & 54 & 50.47 & 0.4620 \\
\bottomrule
\end{tabular}
\end{table}

\begin{figure}[t]
\centering
\includegraphics[width=\linewidth]{figures/fig_budget_sweep.pdf}
\caption{(a) Evasion rate against per-feature budget, $n=\CNsweep$ flows.
Nothing succeeds below $\CEpsFirst$\%. (b) Distribution of the best detector
probability reached at three budgets; the dashed line is $\tau=\CTau$.}
\label{fig:sweep2}
\end{figure}

The curve has a threshold rather than a slope near the origin. Nothing evades at
$5$\% or $10$\%. The first evasions appear at $\CEpsFirst$\%
($\CEvC$ flows, $\CEsrC$\%) and the rate stays there through $20$\%. It reaches
$\CEsrE$\% at $30$\%, and then rises steeply: $\CEsrF$\% at $\CEpsTen$\% and
$\CEsrG$\% at $100$\%.

Two things follow. First, the operational quantity is the perturbation the
feature representation permits, not the sophistication of the search inside it.
A defender who can bound how far a flow's statistics may be pushed while the
attack still functions has bounded the evasion rate, whatever optimiser the
attacker brings. Second, a budget of $50$\% per feature on statistics such as
mean packet length and flow bytes per second is a large perturbation, and
whether an attack surviving it remains the same attack is a problem-space
question this feature-space experiment cannot answer.

%% ===========================================================================
\section{The Reported Rate Is a Property of the Cohort}
\label{sec:cohort}

The pooled ESR of Section~\ref{sec:sweep2} is a weighted average over a cohort
we constructed, and this section shows that the weights matter more than the
optimiser does. Everything below is read off the same \CNsweep{} flows and the
same sweep file; no new run is involved.

\subsection{The attack shifts every score by almost the same amount}
\label{sec:shift}

Start with what the perturbation does to the detector's output rather than with
whether it crosses the threshold. Table~\ref{tab:shift} gives the distribution
of $\Delta p = p_0 - p_\varepsilon$, the drop from the unperturbed probability to
the best probability any arm reaches inside the box. At $\varepsilon = 100$\%
the median drop is \XShiftMed{} with a standard deviation of \XShiftSd{} and an
interquartile range of [\XShiftQone, \XShiftQthree]. Across \CNsweep{} flows spanning
the full range of initial confidence, the attack moves the score by
approximately one fixed amount.

That is not the behaviour of an adaptive search finding per-flow weaknesses. It
is the behaviour of a roughly linear response: within the box the detector's
score is close to an affine function of the ten manipulable features, so pushing
every feature to its permitted extreme buys a displacement that barely depends
on where the flow started. The dispersion that does exist (\XShiftSd{} at
$\varepsilon = 100$\%) is small next to the spread of $p_0$ itself.

\begin{table*}[t]
\caption{Distribution of the score drop $\Delta p = p_0 - p_\varepsilon$ over
$n=\CNsweep$ flows, best arm per flow. The drop is close to a constant at every
budget.}
\label{tab:shift}
\centering
\footnotesize
\setlength{\tabcolsep}{5pt}
\begin{tabular}{@{}lrrlr@{}}
\toprule
$\varepsilon$ & Median & SD & IQR & Max \\
\midrule
$5\%$ & 0.0396 & 0.0295 & [0.0100, 0.0496] & 0.1204 \\
$10\%$ & 0.0396 & 0.0296 & [0.0100, 0.0496] & 0.1204 \\
$15\%$ & 0.0396 & 0.0309 & [0.0100, 0.0600] & 0.1204 \\
$20\%$ & 0.0596 & 0.0278 & [0.0496, 0.0707] & 0.1504 \\
$30\%$ & 0.1207 & 0.0332 & [0.0996, 0.1596] & 0.1918 \\
$50\%$ & 0.2009 & 0.0437 & [0.1700, 0.2526] & 0.2909 \\
$100\%$ & 0.3009 & 0.0392 & [0.2814, 0.3414] & 0.3909 \\
\bottomrule
\end{tabular}
\end{table*}

\subsection{Consequently the starting score, not the attack, decides the outcome}
\label{sec:decides}

If the attack subtracts a near-constant $s$ from every score, then a flow evades
exactly when $p_0 - s < \tau$, that is, when $p_0 < \tau + s$. That prediction
uses nothing about the flow except its unperturbed probability and nothing about
the attack except one median.

Table~\ref{tab:rule} tests it at every budget. At $\varepsilon = 100$\% the rule
``evades if and only if $p_0 < \XThr$'' is correct for \XRuleOk{} of \CNsweep{}
flows, \XRulePct\%, and the point-biserial correlation between $p_0$ and the
evasion indicator is $\XCorr$. The rule is weaker at intermediate budgets, where
the constant shift is small enough that the dispersion around it matters
relative to the distance to the threshold, and it is vacuous at
$\varepsilon \le 10$\% where nothing evades at all. At the budget the headline
number comes from, it is nearly exact.

\begin{table*}[t]
\caption{How much of the outcome the unperturbed probability alone explains. The
predicted threshold is $\tau$ plus the median drop at that budget; the last
column is the point-biserial correlation between $p_0$ and evasion.}
\label{tab:rule}
\centering
\footnotesize
\setlength{\tabcolsep}{5pt}
\begin{tabular}{@{}lrrrr@{}}
\toprule
$\varepsilon$ & Threshold & Correct & \% & $r_{pb}$ \\
\midrule
$5\%$ & 0.5396 & 107/107 & 100.0 & n/a \\
$10\%$ & 0.5396 & 107/107 & 100.0 & n/a \\
$15\%$ & 0.5396 & 105/107 & 98.1 & -0.3863 \\
$20\%$ & 0.5596 & 106/107 & 99.1 & -0.3863 \\
$30\%$ & 0.6207 & 106/107 & 99.1 & -0.6580 \\
$50\%$ & 0.7009 & 73/107 & 68.2 & -0.7240 \\
$100\%$ & 0.8009 & 102/107 & 95.3 & -0.8249 \\
\bottomrule
\end{tabular}
\end{table*}

\subsection{So the pooled rate is whatever the stratum mix makes it}
\label{sec:mix}

Table~\ref{tab:stratumesr} gives ESR by confidence stratum. At
$\varepsilon = 100$\% the three strata report \XEsrBnd\%
(\XEvBnd{} of \XNBnd), \XEsrMed\% (\XEvMed{} of \XNMed) and \XEsrHi\%
(\XEvHi{} of \XNHi). At $\varepsilon = 50$\% they report \XEsrHalfBnd\%,
\XEsrHalfMed\% and \XEsrHalfHi\%. The pooled figure of \CEsrG\% at
$\varepsilon = 100$\% is an average of a rate near one and a rate near zero,
and where it lands is set by how many flows we drew from each band.

The consequence is worth stating arithmetically. Our cohort is
\XNBnd{}/\XNMed{}/\XNHi{} across the three bands. A cohort of the same size
drawn entirely from the high-confidence band would have reported \XEsrHi\%. One
drawn entirely from the medium band would have reported \XEsrMed\%. Both would
have used the identical attack, the identical detector, the identical budget and
the identical code. The span between those two numbers is the whole range that
this experiment can report, and nothing in the attack pipeline determines where
inside it the answer falls.

This is the same lesson as the budget bug of Section~\ref{sec:bug}, arriving
from the other direction. There, a quantity that looked like attack capability
turned out to be an artefact of how the constraint was applied. Here, a quantity
that looks like attack capability turns out to be an artefact of how the
evaluation cohort was assembled. In both cases the number moves by tens of
percentage points without the attack changing at all, and in both cases the
paper reporting it can be entirely correct about its own arithmetic.

\paragraph{What to report instead.} An evasion-rate headline should be
accompanied by the stratification of the cohort over the detector's own output,
and preferably reported per stratum rather than pooled. The per-stratum numbers
are stable statements about the detector --- it is hard to push a flow the
detector is confident about below the threshold, easy for one near it: while
the pooled number is a statement about the sampling frame masquerading as one
about the attack. We report both throughout, and we take the per-stratum table
as the result.

\begin{table*}[t]
\caption{Evasion rate (\%) by confidence stratum and budget, $n=\CNsweep$
flows. Rows are stable properties of the detector; the pooled row is a weighted
average whose weights are our sampling choice.}
\label{tab:stratumesr}
\centering
\footnotesize
\setlength{\tabcolsep}{6pt}
\begin{tabular}{@{}lrrrrrrrr@{}}
\toprule
Stratum & $n$ & $5$\% & $10$\% & $15$\% & $20$\% & $30$\% & $50$\% & $100$\% \\
\midrule
Boundary, $p_0 \in [0.5, 0.6)$ & 7 & 0 & 0 & 29 & 29 & 86 & 100 & 100 \\
Medium, $p_0 \in [0.6, 0.8)$ & 50 & 0 & 0 & 0 & 0 & 0 & 70 & 92 \\
High, $p_0 \ge 0.8$ & 50 & 0 & 0 & 0 & 0 & 0 & 0 & 2 \\
\midrule
Pooled & 107 & 0.0 & 0.0 & 1.9 & 1.9 & 5.6 & 39.3 & 50.5 \\
\bottomrule
\end{tabular}
\end{table*}

\subsection{The oracle over arms buys nothing}
\label{sec:oracle}

One further check, because taking the best of several attacks per flow is a
selection procedure and selection procedures inflate. Table~\ref{tab:armeps}
reports each arm separately against the oracle that takes the minimum over all
three.

The upper-corner arm evades \XArmHiNever{} flows at every budget: pushing the
manipulable features up is simply the wrong direction against this detector.
Greedy coordinate descent matches the oracle at every budget in the sweep
(\XGreedyEqBest), so the oracle's advantage over the single best arm is zero
flows everywhere. The inflation we were checking for does not occur, and we
report the check rather than the reassurance: had it occurred, every ESR in this
paper would need a multiplicity correction, and the reason it does not is that
one arm dominates rather than that best-of-three is a safe construction.

\begin{table*}[t]
\caption{Flows evaded by each arm separately and by the oracle over all three,
out of $n=\CNsweep$. Greedy descent matches the oracle at every budget.}
\label{tab:armeps}
\centering
\footnotesize
\setlength{\tabcolsep}{4pt}
\begin{tabular}{@{}lrrrrrrr@{}}
\toprule
Arm & $5$\% & $10$\% & $15$\% & $20$\% & $30$\% & $50$\% & $100$\% \\
\midrule
Corner, lower & 0 & 0 & 1 & 1 & 6 & 39 & 53 \\
Corner, upper & 0 & 0 & 0 & 0 & 0 & 0 & 0 \\
Greedy coordinate & 0 & 0 & 2 & 2 & 6 & 42 & 54 \\
Oracle best-of-three & 0 & 0 & 2 & 2 & 6 & 42 & 54 \\
\bottomrule
\end{tabular}
\end{table*}

%% ===========================================================================


\section{Discussion}
\label{sec:discussion}

\subsection{A budget bug looks like a capability result}
The failure mode here is worth naming because it is not specific to this
experiment. Any comparison of search strategies under a matched budget is
really two claims: that the strategies differ, and that the budget was equal.
When the budget is enforced incrementally rather than absolutely, the second
claim quietly becomes false in a way that is correlated with the first, because
the strategy that moves most gains the most budget. The result then reads as a
capability difference and is reported as one.

What makes it hard to catch is that every individual step is legal. No single
iteration violates the $\pm\budgetpct$\% rule; only the composition does. A
reviewer checking the perturbation code one step at a time sees nothing wrong.
The check that catches it is not a code review but an assertion on the output:
compare each final vector against the original and confirm every coordinate is
inside the stated box. That assertion costs three lines and would have failed
on the first flow.

\subsection{A one-call baseline should be mandatory}
The corrected table contains a baseline the original did not: evaluating the
detector once at the corner of the box most favourable to the attacker. It
matches or beats twenty iterations of every search strategy, including the
language model.

We recommend it as a standard control in feature-space evasion work. It is one
detector call, it requires no optimiser, and it establishes what the budget
alone yields before any search is credited. A strategy that does not clear it
has not demonstrated search; it has demonstrated the budget. That is precisely
what our earlier result did, and no amount of strategy comparison would have
revealed it, because all four arms were being measured against each other rather
than against the constraint.

\subsection{What a defender should take from this}
The practical reading is narrower than the one we previously offered, and
sturdier. Exposing a confidence score did not, in this setting, let any attacker
cross the threshold within a realistic perturbation budget. The detector's
resistance here comes from the size of the box, not from any property of the
score interface, so the defensive lever is the feature representation: which
statistics an attacker can move, and how far, while the traffic still does what
it was sent to do.

Score hiding remains defensible on other grounds, and we take no position on it
here. What we can no longer offer as evidence for it is this experiment.

%% ===========================================================================
\section{Limitations}
\label{sec:limitations}

\emph{The LLM arm is partial.} \CLlmN\ of \CN\ flows, all from the
high-confidence stratum, because the run ended before completing the cohort.
The bound it supports is $\CLlmCPhi$\%, which is weak. It is enough to say the
model did not evade where the cheap baselines also did not, and not enough to
rank it against them.

\emph{One detector, one attack class, one dataset.} A Random Forest on
CICIDS2017 DDoS flows. A margin-based detector with a smooth score surface may
behave differently under the same search, and the dataset's known labelling and
feature-construction defects are inherited here.

\emph{Feature space, not problem space.} We perturb the flow statistics
directly. Whether a perturbation of the size the sweep requires corresponds to
traffic an attacker could actually emit, and that still achieves the attack, is
not tested and is the harder half of the problem.

\emph{The sweep reports the best arm per flow.} It is therefore an upper bound
on what any one strategy achieves at that budget, chosen to be conservative
about the defender's position rather than favourable to ours.

\emph{The null is a null.} Zero successes in \CN\ trials bounds the rate below
$\CCPhi$\%; it does not establish that the rate is zero. A larger cohort would
tighten the bound and nothing here rules out a rare evasion.

%% ===========================================================================
\section{Conclusion}
\label{sec:conclusion}

We reported that an LLM-guided search evades $\LLMesr$\% of DDoS flows against a
Random-Forest detector while classical optimisers manage almost none. The
per-feature budget in that experiment was enforced against the current vector
rather than the original flow, so the reachable set grew multiplicatively with
the iteration count and the strategy that moved most received the most budget.
The result measured the bug.

With the box fixed to the original vector, every strategy evades $\CGreedyEv$ of
\CN\ flows, and a single detector call at the corner of the box matches twenty
iterations of any search. The language model is no exception: $\CLlmEv$ of
\CLlmN, with the smallest score movement of any arm at the second-highest cost.

What does move the rate is the budget itself, from nothing below $\CEpsFirst$\%
to $\CEsrF$\% at $\CEpsTen$\%. The defensible claim from this work is about how
far a flow may be perturbed, not about who is doing the perturbing.


\section*{Declaration of competing interest}

The authors declare that they have no known competing financial interests or
personal relationships that could have appeared to influence the work reported
in this paper.


\section*{Data availability}

The measured result files, the generator scripts that read them, and the LaTeX
sources that substitute their output into this document are released with the
paper. Every number, table and figure here is produced by running those scripts
against those files; none is transcribed by hand.

The artefact (result files, generator scripts and \LaTeX{} sources) is publicly
available at \url{https://github.com/vithamchoi/nids-evasion-budget-audit}, where the commit that
produced this version is tagged \texttt{v1.0.0}. Re-running the generators in
\path{latex/} against \path{results/} regenerates every table, figure and
inline number in this paper. The same artefact will be deposited in a public
archive under a persistent identifier.


\bibliographystyle{IEEEtran}
\bibliography{refs}

\end{document}
